\chapter{What the banking studies establish}
\label{r3:banking-studies}

The preceding design becomes concrete when the same analysis is followed
from its intervention to its financial account. This chapter uses published
banking designs to explain which information is useful, and reproducible
synthetic studies to test selected estimators and interfaces. The simulated
populations and cash schedules are deliberately specified. None is data
from JPMorgan Chase, Bank of America, HSBC or another institution.

\section{What banking evidence can transfer}
\label{r3:banking-evidence}

Published estimates arrive with an institutional design and population.
In \citet[printed pp.2--4]{drechsler2017}, the deposits-channel argument
uses within-bank branch comparisons across local market concentrations.
Bank-specific lending opportunities are controlled through the comparison;
the interpretation relies on a bank being able to transfer funding across
its branches. The subsequent lending analysis compares banks within a
county and uses their deposit-market exposure. Those are different
aggregation levels and identifying comparisons. The design supports a
mechanism linking market power, deposit pricing and transmission, not
an effect for any new household retention offer.

Relationship banking adds endogenous selection. The mover design in
\citet[Section 4.2.1, equations (6)--(8)]{basten2026} instruments an
existing deposit relationship with the bank's market share in the
household's former municipality, while controlling current share in
the new location. It investigates stickiness and persistent preferences.
The exclusion that prior market availability affects later borrowing
through the relationship remains substantive; the mover population
also limits transfer. This is a useful banking connection to
Chapter~\ref{r3:instruments}, not a universally valid instrument recipe.

For credit supply, the preserved 2006 version of
\citet[Section III.B, equation (5)]{mian2006} compares lenders to the same
firm after differencing. Firm effects remove common borrower-demand
shocks. The paper explicitly discusses the multiple-bank population,
anticipated adjustments and bank-specific demand components that might
remain correlated with funding shocks. A same-borrower comparison is
therefore much stronger than a pooled lending regression, while still
requiring a defensible residual supply variation.

The HARP design of \citet[printed pp.2--3]{agarwal2020} compares program
exposure for similar agency and non-agency high-LTV mortgages. It
distinguishes selected refinancing outcomes from program-level
differences including take-up. The product lesson is to retain eligibility,
assignment, application, acceptance and closing separately, and evaluate
the original eligible population. A borrower cash-flow gain is not
automatically a gain to the bank that owns or services the old mortgage.

Finally, \citet[printed pp.112--113]{beraja2019} explains why regional
refinancing differences do not identify common aggregate responses:
those responses were differenced out. Its model adds collateral and
lender-income channels. \citet[pp.5--6]{benetton2021} combines
discrete-continuous mortgage demand with multiproduct lender pricing,
using capital requirements as cost shifters under additional exclusions.
These sources motivate a structural extension interface when the bank
asks about competitor responses or market-wide regulation. A local
causal estimate may discipline that model but does not determine its
entire equilibrium counterfactual.

\section{A deposit rollout with a complete contribution account}
\label{r3:deposit-study}

The first synthetic study has 96 branches and 30 months indexed from zero.
Twenty-four branches adopt in each of months $12,18,24$, and 24 never
adopt. In USD millions, untreated balances are
\[
 B_{it}(0)=30+\alpha_i+.07t+.4\sin(t/3)+\epsilon_{it},
 \qquad \alpha_i\sim N(0,3^2),
\]
with independent branch disturbances
$\epsilon_{it}=.65\epsilon_{i,t-1}+\eta_{it}$,
$\eta_{it}\sim N(0,.28^2)$. The first disturbance is drawn at the
innovation scale. Cohort labels are independent of these disturbances.
For adoption cohort $g$ and event month $k=t-g\ge0$, the stipulated
balance effect is
\[
 \tau_{gk}=\{.85+.025(g-12)\}(1+.14k).
\]
Untreated parallel trends hold by construction, while treatment effects
differ across cohorts and grow with exposure.

The estimator compares each treated cohort's change from $g-1$ with
the same change in never-treated branches, then averages all treated
cohort-month cells equally. Since cohort sizes are equal, this target
also equally weights treated branch-months. Earlier cohorts receive
more total weight because they contribute more treated months.
The challenger is a single treatment-indicator regression with branch
and calendar effects. Chapter~\ref{r3:panels} explains why its
heterogeneity weights need not equal the requested aggregation.

The financial account uses an annual funding margin $m=.025$, promotional
rate increment $r=.0025$, and monthly program cost $c=1800$ dollars
per treated branch. Let balances now be in dollars. Then
\[
 C(1)=\frac{m-r}{12}B(1)-c,\qquad C(0)=\frac m{12}B(0),
\]
so
\begin{equation}
 \Delta C=\frac{m-r}{12}\Delta B-\frac r{12}B(0)-c.
 \label{r3:deposit-account}
\end{equation}
The second term pays the higher rate on funds that would have remained.
At a USD\,30 million baseline, the additional balance needed to break
even is $(rB(0)+12c)/(m-r)=\$4.293$ million. Counting only the return
on new balances would omit a material cost.

\begin{table}[htbp]\centering\small
\begin{tabular}{@{}lrrrr@{}}
\toprule
Outcome & True effect & Cohort ATT & 95\% interval & TWFE\\
\midrule
Balances (USD m) & 1.780 & 1.689 & [1.564, 1.816] & 1.442\\
Contribution (USD) & -5030.635 & -5218.875 & [-5498.095, -4935.700] & -5607.561\\
\bottomrule
\end{tabular}

\caption{One synthetic deposit rollout. The reference uses the generated
potential outcomes; intervals resample entire branches within cohorts.}
\label{r3:deposit-results}
\end{table}

The cohort estimate is USD\,1.689 million of additional balances and
about USD\,$-5,219$ of monthly contribution per treated branch-month.
The contribution interval is $[-5498,-4936]$. The reference contribution
is about $-5031$. Positive retained balances coexist with a loss after
repricing. The 499-draw branch bootstrap conditions on the cohort design
and common macro path; this one sample is not a coverage qualification.

\begin{figure}[htbp]\centering
\includegraphics[width=.95\textwidth]{derivations/deposit_study.pdf}
\caption{Cohort responses and the signed contribution account under the
specified synthetic law. Different event profiles change aggregation.}
\end{figure}

A targeted offer might lower repricing cost, but it is a new intervention.
Withdrawal risk is not the same as preventable withdrawal. Its evaluation
must use treatment-effect or policy-value information and include the
cost of selection. Repeated offers also alter later balances and
eligibility, invoking the longitudinal model. Internal transfers between
group accounts must be distinguished from external funding retained;
the simple simulation does not model that displacement.

\section{Collections: cumulative histories and policy contrasts}
\label{r3:collections-study}

The collections model has three observed states: current, moderate
arrears and severe arrears. An assistance action is unavailable in the
current state. For no assistance and assistance, respectively, the
row-stochastic transition matrices are
\[
 P_0=\begin{pmatrix}.88&.10&.02\\.30&.50&.20\\.08&.22&.70\end{pmatrix},
 \qquad
 P_1=\begin{pmatrix}.88&.10&.02\\.50&.40&.10\\.18&.42&.40\end{pmatrix}.
\]
The initial state distribution is $(.6,.3,.1)$, and historical assistance
probabilities are $(0,.35,.60)$. Reward after transition is
$(20,-10,-160)$ dollars according to the next state, less $12$ for
assistance; terminal values are $(0,0,-100)$. The horizon is eight
decisions and the discount is $.99$. The unused action-one transition
in the current state is not empirically identified; no evaluated policy
uses it.

Fit transition rows using 3,000 training trajectories and evaluate
three fixed policies on 8,000 independent trajectories. The estimator
uses half-count smoothing on supported transition cells. The policies
give no extra assistance, assist only severe arrears, or assist every
eligible state. Backward dynamic programming computes the reference
values from the generating kernels, which are unavailable to fitting.

\begin{table}[htbp]\centering\small
\begin{tabular}{@{}lrrrrr@{}}
\toprule
Policy & True USD & Direct & IPS & DR & Final ESS\\
\midrule
No extra assistance & -139.58 & -123.44 & -135.08 & -137.56 & 146\\
Severe state only & -58.26 & -49.57 & -59.95 & -59.55 & 1267\\
All eligible & -2.34 & 0.48 & 2.93 & 1.70 & 489\\
\bottomrule
\end{tabular}

\caption{Values in synthetic dollars per trajectory. IPS uses cumulative
ratios; DR adds the continuation residuals. Final-weight ESS is a
diagnostic, not a general effective sample size for every estimator.}
\end{table}

The direct estimate uses fitted transitions alone. Importance sampling
changes the action law along the whole path. The DR score adds discounted
continuation residuals to the initial fitted value. Weighting each
reward only by the current action ratio instead leaves the historical
state distribution in place. In this example that shortcut gives about
$-36.79$ for assistance to all eligible states, compared with the
reference $-2.34$. It answers a different dynamic question.

Policy contrasts must preserve covariance because the same evaluation
customers contribute to each policy. The two prespecified DR contrasts
against no assistance are about $78.01$ and $139.26$; their normal
intervals use a Bonferroni adjustment for two two-sided comparisons,
giving $[44.58,111.44]$ and $[103.83,174.69]$. This is an asymptotic
construction conditional on training, not a finite-sample guarantee or
evaluation after searching over an unrestricted policy class.

\begin{table}[htbp]\centering\small
\begin{tabular}{@{}lrrr@{}}
\toprule
Specification & Bias (USD) & RMSE & 95\% coverage\\
\midrule
DR & 0.15 & 17.33 & 0.938\\
Wrong Q; true logging & 0.71 & 18.28 & 0.869\\
True Q; wrong logging & 0.55 & 10.72 & 0.912\\
Both wrong & 1.41 & 12.67 & 0.856\\
\bottomrule
\end{tabular}

\caption{160 repeated synthetic experiments, each with 600 training and
2,000 evaluation trajectories. Bias identities do not guarantee normal
interval accuracy at this sample size.}
\label{r3:collections-coverage}
\end{table}

Coverage for the fitted DR procedure is $.9375$, with Monte Carlo
standard error about $.019$. With zero continuation regressions and
the true logger it falls to $.86875$, despite the exact population
unbiasedness of that construction. Rare weighted paths make finite-sample
normal approximation fragile. Exact enumeration over a three-decision
version verifies the robustness identities and an independent backward
recursion. It cannot prove reliable eight-decision intervals.
The product should retain these unfavorable results and offer a
better-supported policy or evidence design when the requested history
is poorly represented.

\section{Hidden states: fit the logger before evaluating its replacement}
\label{r3:latent-study}

A separate experiment has two hidden states and three observed signal
categories. The initial hidden distribution is $(.72,.28)$, with
transitions and emission probabilities
\[
 T_0=\begin{pmatrix}.93&.07\\.22&.78\end{pmatrix},\quad
 T_1=\begin{pmatrix}.94&.06\\.52&.48\end{pmatrix},\quad
 B=\begin{pmatrix}.78&.18&.04\\.16&.34&.50\end{pmatrix}.
\]
Actions either have constant probability $.4$, or private-state
probabilities $(.15,.80)$. Fitting receives only signals and actions
from 800 trajectories with 12 transitions. Four-start EM estimates
initial, transition and emission probabilities; the joint model also
estimates private logging probabilities. State labels are ordered by
their mean signal for reporting.

Compare a fixed reactive policy acting only on signal category 2 with
never acting. Its eight-decision reward is next-signal payoff
$(18,0,-90)$ less an $8$-dollar action cost, with discount $.99$.
The generating-model gain is $22.34$.

\begin{table}[htbp]\centering\small
\begin{tabular}{@{}lrrr@{}}
\toprule
Estimation design & Max $P$ error & Max $B$ error & Policy gain (USD)\\
\midrule
Randomized & 0.026 & 0.054 & 24.90\\
Hidden logging, omitted action law & 0.384 & 0.066 & -22.40\\
Hidden logging, joint action law & 0.045 & 0.007 & 13.61\\
\bottomrule
\end{tabular}

\caption{Fitted hidden-state models with unknown emissions. These are
model-conditional estimation comparisons, not a general identification
theorem for observational POMDPs.}
\end{table}

Randomized logging gives a fitted gain of $24.90$. On privately logged
data, omitting the action likelihood gives $-22.40$, while the joint
fit gives $13.61$. The likelihood values of the omitted and joint
models cannot be directly compared as if they were densities of the
same experiment: one omits the action factor. Their relevant comparison
here is the policy target and the specified observational factorization.

The capped optimization must remain visible. In 30 smaller repeated
experiments, all 30 selected omitted-logger fits and seven joint fits
reach the 120-iteration cap. The omitted model's gain bias is about
$-38.95$; the joint model's is $3.41$ with Monte Carlo standard error
$4.02$. These are results for a budgeted fitting procedure. A longer
same-data run, retained separately, checks whether the main sign
reversal persists after its stopping criterion; it is not a replacement
for every capped replicate.
It converges after 163 iterations for the selected omitted-logger fit
and 69 for the selected joint fit, with gains $-22.400$ and $13.612$.
Thus the main sign reversal survives the longer optimization budget.

The randomized-design trajectory bootstrap refits unknown emissions
and transitions in 79 samples and gives a gain interval
$[14.61,38.91]$; one fit reaches its cap. This is a fixed-policy,
model-conditional interval, not uncertainty for a selected optimal
policy or for the causal model class.

Planning on the supplied generating model gives value about $5.303$
on a 501-point belief grid. Doubling to 1,001 points changes it by
about $7.5\times10^{-5}$ dollars. Grid agreement is a numerical
diagnostic, not a certified discretization-error bound. The no-assistance
and fixed reactive values are about $-21.92$ and $.418$.
The optimized plan receives the filtered belief after the current
signal, never a future-data smoother.

\section{Matched lending and shared macro shocks}
\label{r3:credit-dependence-study}

The credit-supply example has 600 firms, 24 banks and three distinct
lenders per firm. Draw bank shock $Z_b$ and firm demand $U_f$
independently, then choose relationships with probabilities proportional
to $\exp(.65U_fZ_b)$. Generated loan growth is
\[
 \Delta L_{fb}=2-3Z_b+4U_f+\eta_b+\epsilon_{fb},
 \quad \operatorname{sd}(\eta_b)=.45,\quad
 \operatorname{sd}(\epsilon_{fb})=1.5.
\]
Selection makes observed lender shocks correlate with firm demand.
Within-firm demeaning removes the common $4U_f$ term. The pooled
coefficient is $-1.080$, while the within-firm coefficient is $-2.923$
against a generating coefficient of $-3$. The illustrative bank-cluster
$t_{23}$ interval is $[-3.027,-2.819]$. This single simulation does
not establish network or few-cluster inference; it isolates the
identifying comparison. A single bank would need outside-lender data
to implement it faithfully.

For dependence alone, generate
$X_{it}=a_i+b_t+\epsilon_{it}$ with independent standard-normal
components, 96 entities and 36 dates. The standard error of the
grand mean is exactly
\[
 \sqrt{1/96+1/36+1/(96\cdot36)}=.19617.
\]
Across 1,200 replications, nominal 95\% interval coverage is about
$.227$ treating rows as independent, $.693$ clustering only by entity,
and $.944$ with entity/date inclusion--exclusion. The last Monte Carlo
standard error is $.0066$. The simulation's negative-variance
truncation rule was never activated; the reusable component instead
returns an explicit flag for indefinite covariance.

The structural-series example uses
$F=\left(\begin{smallmatrix}.7&0\\.15&.6\end{smallmatrix}\right)$
and reduced innovation covariance
$\left(\begin{smallmatrix}1&.8\\.8&1\end{smallmatrix}\right)$.
Two recursive orderings have the same reduced law but normalized
funding-shock impacts $(1,.8)'$ and $(1,0)'$. Their lending responses
differ. The observational fit cannot determine which timing restriction
is economically appropriate. This complements the larger-sample
variance comparison: more accurate uncertainty about an observed law
does not eliminate structural ambiguity.

\section{Joint response inference, persistence and learned adjustment}
\label{s4:response-results}
The response experiments instantiate Section~\ref{s4:response-construction}
with two banking outcomes, 13 horizons from zero through twelve, and a
common set of origins. Each row of Table~\ref{s4:response-coverage} contains
500 independent replications. Whole-curve coverage requires all 26 true
coordinates to lie in their reported sets. Applying 95\% pointwise intervals
to every plotted horizon is therefore deliberately an adverse comparison.
Zero numerical failures occurred; failed replications, had they occurred,
would remain in the denominator.

\begin{table}[htbp]\centering\small
\begin{tabular}{lrrrrrr}
\toprule
Case & Iid & HAC & Joint & Max & IV set & VAR\\
\midrule
Strong, white & 0.434 & 0.414 & 0.958 & 0.930 & 0.952 & 0.950\\
Strong, persistent & 0.132 & 0.452 & 0.934 & 0.886 & 0.930 & 0.938\\
Persistent, long & 0.116 & 0.468 & 0.926 & 0.876 & 0.928 & 0.978\\
Weak, persistent & 0.670 & 0.850 & 0.962 & 0.944 & 0.924 & 0.986\\
Near unit & 0.246 & 0.598 & 0.818 & 0.724 & 0.776 & 0.966\\
Invalid IV & 0.002 & 0.012 & 0.182 & 0.136 & 0.270 & 0.000\\
\bottomrule
\end{tabular}

\caption{Whole-curve coverage of nominal 95\% procedures. ``Joint'' uses
Bonferroni; ``max'' uses the estimated Gaussian maximum; ``IV set'' inverts
the joint residual tests. Standard Monte Carlo errors near .95 are about
.010. These are synthetic experiments, not coverage guarantees.}
\label{s4:response-coverage}
\end{table}

Strong instruments have innovation loading one; the weak case has loading
$.06$. The reference system has spectral radius $.81$; the near-unit case
uses $.985$. Persistent instrument noise has AR coefficient $.85$, rising
to $.95$ in the near-unit case. Most experiments use 800 origins; the long
persistent case uses 5,000. The invalid instrument additionally loads on
another structural innovation by $.6$. The initial stationary law and
all generating matrices are fixed in the reproducible specification.

Several method decisions follow. Joint covariance matters because a
pointwise interval is not a curve guarantee. Yet a consistent HAC estimator
can still underestimate finite-sample uncertainty: the Gaussian-maximum
band covers only .886 and .876 in the two persistent strong-instrument
experiments, and .724 near the unit boundary. Bonferroni is more conservative,
but its .818 near-unit coverage also fails the preset adequacy threshold.
These failures prohibit presenting the default bandwidth as automatically
calibrated for persistent bank series. The statistical result remains
asymptotic under a fixed stability margin, not uniform near a unit root.

The proxy VAR has much smaller mean squared error in the correctly
specified finite-order generating law. That advantage is a consequence of
its propagation restriction. It does not establish that an economist should
replace direct moments when finite-order dynamics are misspecified. Weak
normalization makes the direct and proxy mean squared errors about 752
and 13.5; 96.4\% of runs yield an unbounded residual-test set in at least
one coordinate. A narrow decision recommendation would conceal the
information deficiency. Instrument invalidity defeats both routes, including
zero proxy-VAR coverage in this example.

A separately specified follow-up, with a new random seed, compares bandwidth
$L=\lfloor n^{1/4}\rfloor$ and $4L$ at 5,000 origins. Their whole-curve
coverages are .948 and .968, with mean widths .312 and .338. This new
experiment does not replace the original .926 finding or establish that
$4L$ is universally preferable. The selected workflow reports sensitivity
to the lag window, compares the structural challenger when justified and
withholds a calibrated finite-sample claim in regimes that fail the stated
qualification. Picking the more favorable seed or interval after observing
the answer would be a different, unqualified procedure.

The integrated experiment has one common random-walk trend and two stationary
spreads. Mapping stationary-change IV estimates to cumulative level responses
gives whole-curve coverage .966 and mean squared error .0479. Treating the
levels as stationary gives .586 coverage and error 55.66. Cointegration is
substantive here: it defines which components are stationary and which
intervention response is being summed. The comparison does not certify
unknown rank selection or arbitrary local-to-unit triangular arrays.

With 4,000 dates and 80 bounded controls, separated-block sparse IV gives
.942 coverage of the scalar effect $1.2$, compared with .950 for oracle
nuisances. A three-factor reduction and an unadjusted estimator both have
zero coverage, with mean squared errors about .058 versus .00102 for
sparse adjustment. The causal confounder is sparse in this generator and
need not align with its leading covariance components. Factor compression
requires an information-preservation argument, not only a prediction score.
The fixed penalty multiplier in this experiment establishes feasibility;
the required population nuisance rates remain assumptions to justify in
each application.

For a fixed 96-bank panel, the additional mean-score experiment separates
common and bank-specific Gaussian AR components. At 800 dates, row-iid
intervals cover in .116 of runs; date-aggregate HAC covers in .856; an
oracle using the exact finite-sample Gaussian variance covers in .938.
At 2,000 dates the corresponding values are .104, .890 and .946.
Correctly recognizing the date as the information unit removes a gross
variance error, but the short Bartlett lag window still underestimates
persistent uncertainty. More customer rows do not repair either problem.

\section{Complete computational workloads and their statistical meaning}
\label{s4:workflow-results}
Table~\ref{s4:workflow-cost} measures fresh processes, including imports,
synthetic generation, local Parquet write/read, observation reconstruction,
fitting, uncertainty, the specified comparison or resampling, and report
serialization. Stage timings, physical and logical I/O, peak resident
memory, seeds, numerical differences and optimizer stops are retained.
Each row is the median of three processes on an AMD EPYC 7773X host with
256 logical CPUs; BLAS, DuckDB and Polars are each restricted to one thread.
The machine and library versions are recorded with the results. The
measurements are local warm-filesystem workloads, not network/object-store
or distributed-system benchmarks.

\begin{table}[htbp]\centering\small
\begin{tabular}{lrrr}
\toprule
Workload & Raw rows & Seconds & Peak MiB\\
\midrule
Response, small & 77,952 & 2.00 & 267.7\\
Response, medium & 772,608 & 2.22 & 356.8\\
Response, large & 3,849,216 & 3.38 & 717.2\\
Learning, small & 4,000 & 1.56 & 256.0\\
Learning, large & 24,000 & 2.07 & 432.0\\
Latent, small & 3,200 & 7.49 & 254.0\\
Latent, large & 38,400 & 48.27 & 264.0\\
Temporal graph, small & 1,000 & 1.49 & 250.0\\
Temporal graph, large & 4,000 & 1.71 & 257.2\\
\bottomrule
\end{tabular}

\caption{Complete synthetic workload cost. Response rows include three raw
versions per bank/date; learning rows include two versions per date.
Latent rows include 19 whole-trajectory bootstrap refits with three starts.
Memory includes the interpreter, libraries and retained synthetic arrays.}
\label{s4:workflow-cost}
\end{table}

For the response workload, bank deviations are centered across banks so
that the fixed weighted aggregate equals the common structural series.
Raw data contain an initial publication, an eligible correction and a
future revision. The SQL operation restricts availability before selecting
the highest eligible version, then aggregates to dates. It is checked
against the separately retained generating series, and the fitted curve
against its direct numerical reference. Maximum discrepancies are below
$10^{-9}$ at every scale. The largest case reconstructs 3,849,216 versioned
records for 256 banks and 5,012 dates, then fits direct and proxy responses,
joint covariance, residual-test sets and a bandwidth comparison. It is
one common-shock analysis; the record count does not create millions of
independent macroeconomic shocks.

The learned-control workload reconstructs a wide versioned table and fits
both separated-block sparse adjustment and the factor challenger, retaining
their errors and covariance. The large case has 12,000 dates and 200 controls.
The temporal graph workload fits 12 variables at 4,000 dates and propagates
the joint coefficient box through 13 horizons. Its timing does not establish
uniform high-dimensional graph coverage; the inferential reference remains
the declared fixed-dimensional structural class. Exponential DAG enumeration
and the 25-DAG Gaussian reference are different workloads.

The latent workload reconstructs complete customer trajectories, fits the
unknown-emission model with multiple starts, computes observed-information
value uncertainty and refits every bootstrap sample. The larger case has
3,200 customers and 12 decisions. Its model fit and information calculation
cost only a small fraction of the 19-refit workflow. Nineteen bootstrap
draws are suitable for demonstrating cost scaling, not precise 95\% interval
endpoints. The separate 79-refit fixed-policy example and repeated-sampling
study address different inferential questions. All caps, boundary and
optimizer diagnostics remain in the receipts, including intermediate EM
caps followed by likelihood refinement.

These results change the engineering priority. For short response curves,
data reconstruction and process overhead can dominate the numerical moment
calculation. For controlled likelihoods, repeated fitting dominates.
Caching immutable views helps the first workload; efficient likelihood
scores, warm initial proposals checked against independent starts and
parallel independent bootstrap fits can help the second. Each proposed
acceleration must retain the same target and numerical tolerance. No
bank-level labor saving follows from these CPU timings: the expert
composition comparator can use the same engines and estimators.


\section{The implemented data and score components}
\label{r3:component-evidence}

Two small local components implement selected parts of
Chapter~\ref{r3:construction}. The temporal component constructs
Equation~\eqref{r3:revision-feature} under its declared inclusive
availability convention. In an adversarial fixture, the correct
early and later amounts are $160$ and $90$. Using the latest
retrospective values gives $90$ at the early decision; filtering
economic time before choosing the revision gives $110$ at the later
decision. Delayed data, voids and a moved economic date explain the
differences.

The component rejects duplicate event revisions, incompatible currencies,
unqualified feed completeness and naive timestamps. Reconstruction review
also added explicit rejection of empty or non-string event/decision
identities. Without it, a missing decision entity could return a
misleading zero. This is a per-decision semantic reference, not a
bank-scale temporal query engine.

The score reducer retains raw sums/counts for entities, dates and their
intersections, centers globally and computes the stated crossed
covariance. Independent batch calculations agree to numerical precision
for both unique and repeated entity/date cells. Partition merges preserve
the result. Invalid multidimensional cluster keys now fail before
changing accumulator state. The caller still owns the target weighting,
Jacobian, training split and dependence argument; duplicate-task
protection remains part of the proposed execution layer.

\begin{table}[htbp]\centering\small
\begin{tabular}{@{}rrrr@{}}
\toprule Rows & Total seconds & Reduce seconds & Peak RSS (MiB)\\\midrule
100,000 & 0.097 & 0.078 & 77.1\\
1,000,000 & 0.233 & 0.133 & 115.2\\
10,000,000 & 1.186 & 0.630 & 123.1\\
\bottomrule\end{tabular}

\caption{Median of three fresh processes at each size on the recorded host.
The workload generates and reduces precomputed two-dimensional scores;
it excludes source I/O, history reconstruction and model fitting.}
\label{r3:scale-measurement}
\end{table}

The largest run uses ten million scores, 10,000 declared entities and
120 dates. Its dense intersection allocation grows with the product
of those group counts. Peak memory includes the interpreter and
libraries. This measurement supports only the bounded reduction
operation; it is not ten million fitted causal analyses or evidence
of distributed end-to-end bank performance.

The combined evidence establishes several useful construction tests:
incorrect cash accounts can reverse a business conclusion, current-only
weights can answer the wrong dynamic question, omitted private logging
can reverse a policy estimate, and row counts can conceal shared shocks.
The additional constructions above qualify local study recipes and selected
package operations. Complete product workspaces, institutional validation
and accepted sibling adoption remain distinct engineering and review tasks.

\section{A public household-finance analysis with unknown causal truth}
\label{s5:public-finance}
Public data expose choices that a known-truth simulation cannot. The archived
401(k) application extract contains 9,915 observations from the SIPP 1991
design used in the DML literature \citep[Section 6.2]{dml2018,public401k}.
The outcome is net financial assets; eligibility and participation are
separate binary variables. The declared controls are age, income, education,
family size, marriage, two-earner status, defined-benefit pension status,
IRA participation and home ownership. All required fields are observed in
the retrieved extract, so no row is removed by the declared missingness rule.
The dollar values are those in the archive, without an inflation conversion.

The first question is the average effect of eligibility on assets under
conditional exchangeability given these controls. This is an identifying
assumption about employer/program selection and saving behavior, not a
consequence of fitting a flexible model. The second question is participation's
complier effect with eligibility as an instrument. It adds exclusion,
instrument independence and monotonicity. The ratio of the adjusted
eligibility effect on assets to its adjusted effect on participation yields
that target under those restrictions. Ordinary participation regression
would answer neither question without a different selection argument.
The graph and observation choices must also establish which controls are
appropriate pre-exposure covariates; the cross-sectional extract alone cannot
verify their timing or rule out unmeasured saving preference.

Two nuisance families are compared: standardized quadratic features with
ridge outcome regression and logistic propensity regression, and histogram
gradient boosting for both roles. Three folds and seeds 7101--7103 are fixed
in advance. The original paper's 100-split aggregation and learner collection
are not replicated numerically; this reproduces the design and score with
declared alternative learners. The matched DoubleML check isolates score
agreement from those learner differences.

\begin{center}\small
\begin{tabular}{llrr}\toprule
Nuisance family&Seed&Eligibility ATE&Participation LATE\\\midrule
Quadratic ridge/logit&7101&8,223&11,905\\
&7102&8,391&12,089\\
&7103&8,932&12,935\\
Histogram boosting&7101&7,968&11,398\\
&7102&8,502&12,261\\
&7103&8,299&11,879\\\bottomrule
\end{tabular}\end{center}
These estimates use archived dollars, with the LATE interpreted per
complier under its additional assumptions. For the first quadratic run,
the eligibility 95\% interval is approximately $[5,737,10,710]$; for the
first boosting run it is $[5,683,10,253]$. First-stage estimates range from
.691 to .699. The three quadratic fits each clip one propensity; the
boosting fits clip none. The largest realized inverse-action weight ranges
from about 32 to 47. Small clipping counts do not prove adequate population
overlap. The raw eligibility mean difference is about 19,559, substantially
larger than the adjusted estimates, but that comparison does not reveal the
unknown causal truth.

The reported covariance uses an iid household working law. The extract does
not supply survey weights, strata or primary sampling units. Consequently
these intervals are a reproducible working analysis, not a design-valid
survey inference claim. Alternative control timing, employer selection,
survey design and weak support require sensitivity or additional data.
The result does not transfer automatically to a bank offer, a different
country or current wealth. The useful product behavior is to preserve the
eligibility/participation target distinction, identical matched-reference
inputs, overlap diagnostics, empirical uncertainty and unverified design
assumptions through the final result.


\section{Public monetary histories and the boundary of an empirical claim}
\label{s6:public-dynamic}
The archived statsmodels macroeconomic extract contains 203 quarterly
observations from 1959Q1 through 2009Q3, assembled from FRED and labor
statistics with source access dated 15 December 2009
\citep{public_macrodata_dynamic}. It includes an end-quarter M1 stock,
a quarterly average of monthly three-month Treasury-bill rates, output
and inflation. Its data, source notes and exact numerical implementation
are preserved. This is a real dynamic financial-data analysis with a known
historical snapshot; it is not a current bank deposit-offer dataset.

The fixed projection response is
$100(\log M1_{t+h}-\log M1_{t-1})$ per percentage-point change in
the quarterly-average bill rate. Two lags each of the rate change, money
growth, output growth and inflation enter every regression. Four
specifications are fixed before fitting: these lagged controls; adding
current output growth and inflation; a pre-1984 origin window; and a
post-1984 origin window. Horizons zero through four use common origins
within each specification. Bartlett lag four covariance is computed from
the complete OLS influence, including controls. An independent statsmodels
OLS/HAC call agrees in coefficient and covariance to below $3\times10^{-14}$.

At horizon four, the full-sample lag-control estimate is $-2.102$
percentage points of cumulative money growth, standard error .452.
Adding current output/inflation gives $-2.163$ (.453). The pre-1984
estimate is $-.997$ (.262), while the later window gives $-4.329$ (.685).
There are 196 common full-sample origins and 97/99 in the two windows.
The different windows define different projection populations; selecting
one after observing its significance would add a selection problem. Their
contrast is a registered sensitivity analysis, not proof of a known break
date or a homogeneous policy effect.

These numbers do not identify a bank's price intervention. A quarterly
average rate and end-quarter stock do not supply a within-quarter decision
clock; central-bank information and money demand can confound the
innovation; M1 is not one bank's retained deposits; structural and
measurement regimes can change over fifty years. Current output controls
may be post-policy variables under one causal ordering and predecision
information under another. The data alone cannot settle that ordering.
The inference reported is conditional projection uncertainty, not uncertainty
over all these causal interpretations.

The public-data investigation also inspects the Federal Reserve H.8
description \citep{fed_h8_dynamic}. Its aggregate bank balance-sheet
statistics are useful for surveillance, but do not reveal customer offers,
the private assignment information, first-available entity-linked histories
and representative audited matches required by the treasury brief. Neither
that feed nor the archived macro series identifies the stipulated bank
experiment. Obtaining a current bank-specific price/offer history and its
assignment evidence is therefore a genuine empirical requirement, rather
than an excuse to call any public time series causal.

The central known-truth evidence uses the realistic but explicitly synthetic
funding mechanism, common shocks and revision/linkage construction. At 360
and 1,440 dates, 400 repetitions give financial interval coverage .945 and
.9375; pretending the branches supply independent calendar information
gives .230 and .2175. The corresponding joint Bonferroni response coverage
is .9425 and .965. Binomial Monte Carlo error is approximately one percentage
point for coverage near .95 at 400 repetitions; these differences do not
certify exact nominal calibration. The pair-linkage and audited-servicing
experiments add separate observation and information failures. Taken
together, the public analysis challenges economic interpretation, while
the synthetic analysis checks statistical behavior when its assumptions
are known. Neither estimates a named banking group's causal benefit.

